\subsection{LASSO and Elastic Net}
\label{sec:sparse_methods}
We used logistic-regression versions of LASSO and Elastic Net as sparse predictive baselines \cite{tibshirani1996lasso,zou2005elastic}. Both minimise a class-weighted logistic loss with a penalty on the gene coefficients $\boldsymbol{\beta}$ of the form
\begin{equation}
\Omega(\boldsymbol{\beta})=\alpha\lVert\boldsymbol{\beta}\rVert_1+
\frac{1-\alpha}{2}\lVert\boldsymbol{\beta}\rVert_2^2.
\label{eq:sparse_penalty}
\end{equation}
LASSO uses $\alpha=1$, whereas Elastic Net mixes $L_1$ and $L_2$ regularisation. The final LASSO models were fitted with \texttt{liblinear} coordinate descent \cite{fan2008liblinear}, using balanced class weights, tolerance $10^{-6}$ and inner-cross-validation selection over $C\in\{0.001,0.01,0.1,1\}$. Elastic Net was fitted with R \texttt{glmnet} \cite{friedman2010glmnet} at convergence threshold $10^{-12}$ followed by active-set Newton refinement. For both models, optimality was checked against Karush--Kuhn--Tucker conditions, including inactive coefficients and the intercept. Coefficients with absolute value greater than $10^{-8}$ defined the selected panel.

For each outer LOCO split, $(C,\alpha)$ was chosen by nested inner cross-validation within the three development cohorts using mean validation AUROC, with ties favouring stronger regularisation. Preprocessing (CPM $>1$ in at least 10\% of training samples, nonzero variance, $\log_2(1+\mathrm{CPM})$ and standardisation) was refitted strictly inside every training split; genes with $|\beta|>10^{-8}$ formed the selected signature. To evaluate gene agreement across patient-disjoint cohorts (zero shared samples), we also fitted LASSO and Elastic Net separately within each individual cohort using stratified inner cross-validation; their training-set predictions are not reported as external test results.
